\section{$2$-dynamic coloring of $C_p(1,3)$}

We use color words cyclically on the vertices
$0,1,\ldots,p-1$; as usual, $W^q$ denotes the concatenation of
$q$ copies of a word $W$.

\begin{theorem}\label{thm:circulant-r2}
Let $p\ge6$. Then
\[
\chi_2^d(C_p(1,3))=
\begin{cases}
3, & p\equiv0\pmod9,\\
4, & \text{otherwise}.
\end{cases}
\]
\end{theorem}

\begin{proof}
If $9\mid p$, the word
$
(012120201)^{p/9}
$
defines a $2$-dynamic $3$-coloring: each vertex has a color different
from its neighbors at distances $1$ and $3$, and its neighborhood
contains both other colors. Hence
$\chi_2^d(C_p(1,3))=3$ by
Equation~\ref{eq:absolute-bounds}.

We next give a $2$-dynamic $4$-coloring for every $p\ge6$.
Let $A=0123$ and write $p=4q+r$, $0\le r\le3$. Use the following
cyclic color words:
\[
\begin{array}{c|cccc}
r&0&1&2&3\\ \hline
\phi&A^q&A^{q-1}03231&A^q01&A^q201
\end{array}
\]
A direct verification shows that $\phi$ is a $2$-dynamic coloring.
Thus $\chi_2^d(C_p(1,3))\le4$ for every $p\ge6$.

It remains to prove that a $2$-dynamic $3$-coloring can exist only if
$9\mid p$. We may assume $p\ge7$, since
$C_6(1,3)\cong K_{3,3}$ has no $2$-dynamic $3$-coloring.
Let $\phi$ be a $2$-dynamic $3$-coloring and define
\[
s_i=
\begin{cases}
1,&\phi(i-1)\neq\phi(i+1),\\
0,&\phi(i-1)=\phi(i+1).
\end{cases}
\]
We first observe that $s_i=s_{i+1}=1$ is impossible. Indeed, in this
case both
$
\phi(i-1),\phi(i),\phi(i+1)
$
and
$
\phi(i),\phi(i+1),\phi(i+2)
$
contain the three colors. Hence $\phi(i-1)=\phi(i+2)$, contradicting
the fact that $i-1$ and $i+2$ are adjacent.

Moreover,
$
s_{i-2}=s_i=s_{i+2}=0
$
is impossible, since it would imply
$
\phi(i-3)=\phi(i-1)=\phi(i+1)=\phi(i+3),
$
contrary to the $2$-dynamic condition at $i$.

Since not all $s_i$ can be zero, choose $i$ with $s_i=1$.
Then $s_{i-1}=s_{i+1}=0$, and the previous observation forces
$s_{i-3}=s_{i+3}=1$. By the first observation,
$s_{i-2}=s_{i+2}=0$.
Reconstructing $\phi$ yields $\phi(i-4)=\phi(i+5)$.
Repeating the same argument gives
\[
\phi(i-4)=\phi(i-4+9k)
\]
for every $k$, and hence $p\equiv0\pmod9$.
\end{proof}